\section{Step 27 Write height and its LF}\label{ppm-s27}

Apply the same quotient/remainder method to height, then store LF after its last digit. Keep track of where the height token began: the next region writes Maxval relative to that point before advancing the cursor across the whole suffix.

Author region \texttt{e\_height}, keeping its BEGIN/END markers and the supplied caller, checkpoint labels and cleanup paths. Stop at \texttt{after\_e\_height} for the observation.

Use recipe \texttt{s27-1} with the cumulative source through this edit.

\subsection{Expected observation}
The height and following LF are present; Maxval positions remain sentinel bytes.

\Needspace{9\baselineskip}
\subsection{Check your understanding}
Why test heights nine and ten?

\begin{itemize}
  \item \textbf{A} They change the sample width.
  \item \textbf{B} They exercise one-digit and two-digit token lengths.
  \item \textbf{C} They change the file magic.
\end{itemize}
